\title{DeepSeekMath: Pushing the Limits of Mathematical Reasoning in Open Language Models}

\begin{abstract}
Mathematical reasoning remains a formidable obstacle for language models because of its inherently intricate and structured requirements. In this study, we present DeepSeekMath~7B, developed by further pre-training DeepSeek-Coder-Base-v1.5 7B on a corpus of 120B math-focused tokens gathered from Common Crawl, along with accompanying datasets in natural language and code. DeepSeekMath~7B attains a strong score of 51.7\% on the competitive MATH benchmark, achieving this result without dependence on external toolchains or ensemble methods, and approaching the performance levels exhibited by models such as Gemini-Ultra and GPT-4. Moreover, employing self-consistency sampling across 64 instances enables DeepSeekMath~7B to reach 60.9\% accuracy on MATH. The advancements in mathematical reasoning capability can be ascribed to two central innovations: (1) We unlock the vast potential of publicly available online resources using a rigorously crafted data filtering and selection workflow; (2) We propose Group Relative Policy Optimization (GRPO), an adaptation of Proximal Policy Optimization (PPO), which not only boosts mathematical reasoning proficiency but also brings improvements to PPO’s memory efficiency.
\end{abstract}

\section{Introduction}

Large language models (LLMs) have transformed methods of mathematical reasoning within artificial intelligence, driving notable progress in benchmarks for both quantitative reasoning \citep{MATH} and geometry reasoning \citep{trinh2024solving}. Additionally, these models have played a pivotal role in supporting humans with the resolution of sophisticated mathematical tasks \citep{tao}. Despite these advances, state-of-the-art systems like GPT-4 \citep{gpt4} and Gemini-Ultra \citep{gemini} remain closed to the public, and open-source alternatives currently lag significantly in terms of their capabilities.

This paper presents DeepSeekMath, a specialized language model for mathematical content that surpasses existing open-source alternatives in mathematical reasoning and achieves near-parity with GPT-4 on several academic tasks. To facilitate this advancement, we constructed the DeepSeekMath Corpus—a high-quality pre-training dataset containing 120B mathematics-related tokens. Data collection involved leveraging a fastText-based classifier \citep{joulin2016fasttext} to filter relevant content from Common Crawl (CC). For initial classifier training, positive samples were drawn from OpenWebMath \citep{openwebmath}, while negative samples consisted of assorted unrelated web content. After the first training round, the classifier was utilized to extract more candidate math-related content from CC, which subsequently underwent further filtration and annotation by human experts. This newly curated data was used to further retrain and enhance the classifier’s discriminatory capabilities. Evaluations of our DeepSeekMath-Base 7B model demonstrate the superior quality of the resulting corpus, as evidenced by a 64.2\% accuracy on GSM8K \citep{gsm8k} and a 36.2\% score on the MATH benchmark \citep{MATH}, outperforming models such as Minerva 540B \citep{minerva}. Due to the multilingual nature of the DeepSeekMath Corpus, we also observe gains on Chinese-language mathematical tasks \citep{wei2023cmath,agieval}. We suggest that our methodology for mining and curating mathematical text provides a foundation for further advances, leaving ample space for ongoing community improvement.

DeepSeekMath-Base is built upon the DeepSeek-Coder-Base-v1.5 7B \citep{deepseek-coder}, as our findings suggest that initializing with a model trained on code yields superior results relative to commencing with a conventional general-purpose LLM. In addition, we find that training on mathematical data augments the model’s performance not only in mathematics-specific tasks but also on the MMLU \citep{mmlu} and BBH \citep{bbh} benchmarks. This improvement suggests that mathematical pretraining boosts both mathematical proficiency and broader reasoning skills.

Following pre-training, we conduct mathematical instruction tuning on DeepSeekMath-Base, utilizing data comprising chain-of-thought \citep{cot}, program-of-thought \citep{pot,pal}, and tool-integrated reasoning \citep{tora}. The finalized model, DeepSeekMath-Instruct 7B, surpasses all other 7B-scale models and demonstrates performance on par with open-source instruction-tuned models of 70B parameters.

In addition, we propose Group Relative Policy Optimization (GRPO), a modified reinforcement learning (RL) approach derived from Proximal Policy Optimization (PPO) \citep{schulman2017proximal}. Unlike traditional PPO, GRPO dispenses with the critic model by inferring the baseline using group scores, thereby considerably lowering the computational requirements for training. With only a fraction of English instruction tuning examples, GRPO achieves notable advancements beyond the robust DeepSeekMath-Instruct on both in-domain datasets (GSM8K: 82.9\% $\rightarrow$ 88.2\%, MATH: 46.8\% $\rightarrow$ 51.7\%) and out-of-domain mathematical benchmarks (e.g., CMATH: 84.6\% $\rightarrow$ 88.8\%) during the RL stage. We further introduce a unified framework to contextualize approaches such as Rejection Sampling Fine-Tuning (RFT) \citep{yuan2023scaling}, Direct Preference Optimization (DPO) \citep{dpo}, PPO, and GRPO. Through this integrated perspective, we demonstrate that each of these strategies can be interpreted as forms of direct or streamlined RL algorithms. Our work includes a comprehensive suite of empirical analyses—such as comparisons between online and offline RL, outcome versus process-based feedback, and evaluations of single-turn and iterative RL—to rigorously examine the foundational aspects of this paradigm. Finally, we elucidate how our RL methodology enhances instruction-tuned model performance and discuss future directions to further refine RL efficacy within this unified paradigm.

\subsection{Contributions}
We present a scalable approach to math pre-training, as well as a comprehensive investigation and assessment of reinforcement learning methodologies.

\textbf{Math Pre-Training at Scale}

\begin{itemize}[topsep=0pt]
    \item In this study, we demonstrate that the openly available Common Crawl dataset offers significant resources for mathematical research. Through the use of a carefully crafted data selection pipeline, we assemble the DeepSeekMath Corpus—a superior quality dataset encompassing 120B tokens derived from web pages curated for mathematical relevance. This collection is nearly sevenfold larger than the math-specific web pages utilized by Minerva \citep{minerva} and approximately nine times bigger than the newly published OpenWebMath dataset \citep{openwebmath}.

\item DeepSeekMath-Base 7B, our pre-trained foundational model, demonstrates similar performance to Minerva 540B \citep{minerva}, suggesting that sheer parameter count is not the sole determinant of effectiveness in mathematical reasoning tasks. High-quality pre-training data enables smaller models to attain competitive results.

\item We present the results of our experiments involving math training. Conducting code training before math training enhances the models' performance in mathematical problem-solving, regardless of whether tools are employed. This provides some evidence toward resolving the enduring query: \textit{does code training improve reasoning abilities?} Our results suggest that it does, particularly in the context of mathematical reasoning.

\item While it is a standard practice to utilize arXiv papers for training, particularly within mathematics-focused research, such data did not yield any significant enhancements on any of the mathematical benchmarks used in this study.
\end{itemize}

\textbf{Exploration and Analysis of Reinforcement Learning}

\begin{itemize}[topsep=0pt]
    \item We propose Group Relative Policy Optimization (GRPO), a resource-efficient reinforcement learning technique. Unlike conventional methods that rely on a critic model, GRPO determines the baseline using group scores, which considerably lowers computational requirements relative to Proximal Policy Optimization (PPO).
    \item We show that applying GRPO notably boosts the instruction-tuned capabilities of DeepSeekMath-Instruct when trained only with instruction-tuning datasets. Additionally, reinforcement learning with GRPO leads to improvements in out-of-domain generalization.
    \item We offer a cohesive framework that contextualizes various approaches, including RFT, DPO, PPO, and GRPO. Our study involves comprehensive experimentation on topics such as online versus offline learning, process versus outcome supervision, and comparisons between single-turn and iterative reinforcement learning, allowing for a deeper analysis of the paradigm’s crucial aspects.
    \item Leveraging our unified framework, we investigate the underlying factors contributing to the success of reinforcement learning and outline several promising avenues for enhancing the reinforcement learning of LLMs.
\end{itemize}

\subsection{Summary of Evaluations and Metrics}

\begin{itemize}[topsep=0pt]
    \item \textbf{English and Chinese Mathematical Reasoning}: We systematically evaluate our models using a range of benchmarks in both English and Chinese, spanning mathematical questions from elementary through collegiate complexity. For English, benchmarks assessed include GSM8K \citep{gsm8k}, MATH \citep{MATH}, SAT \citep{llemma}, OCW Courses \citep{minerva}, and MMLU-STEM \citep{mmlu}. The Chinese datasets evaluated comprise MGSM-zh \citep{mgsm}, CMATH \citep{wei2023cmath}, Gaokao-MathCloze \citep{agieval}, and Gaokao-MathQA \citep{agieval}. We assess both the proficiency of the models in producing comprehensive, self-contained textual solutions independently and their effectiveness when leveraging Python-based problem-solving.

Across English-language benchmarks, DeepSeekMath-Base performs on par with the proprietary Minerva 540B model \citep{minerva}, and outperforms all open-source base models—including Mistral 7B \citep{mistral} and Llemma-34B \citep{llemma}—regardless of whether these models have been subjected to mathematical pre-training, frequently by a large margin. In particular, DeepSeekMath-Base achieves even greater success on Chinese-language evaluations, likely due to our decision to diverge from prior approaches \citep{minerva,llemma} that restrict mathematical pre-training to English datasets, instead incorporating a substantial quantity of high-quality non-English data. Following the application of mathematical instruction tuning and reinforcement learning, DeepSeekMath-Instruct and DeepSeekMath-RL achieve impressive results, surpassing 50\% accuracy on the competition-grade MATH dataset—marking a first for open-source models in this area.

\item \textbf{Formal Mathematics}: We assess DeepSeekMath-Base on the informal-to-formal theorem proving benchmark, following the methodology introduced in \citep{dsp_proof}, employing the miniF2F \citep{minif2f} suite and selecting Isabelle \citep{isabelle} as the proof assistant. In this context, DeepSeekMath-Base exhibits notable few-shot capabilities for autoformalization.
\item \textbf{Natural Language Understanding, Reasoning, and Code}: To gain a holistic understanding of the model's proficiency in language comprehension, reasoning, and programming, we subject DeepSeekMath-Base to evaluation with the Massive Multitask Language Understanding (MMLU) benchmark \citep{mmlu}—which comprises 57 diverse, multiple-choice subject areas—alongside BIG-Bench Hard (BBH) \citep{bbh}, featuring 23 demanding problems that predominantly necessitate multi-step logical reasoning. Furthermore, we test its coding aptitude through the HumanEval \citep{codex} and MBPP \citep{mbpp} datasets, both of which are established benchmarks for assessing programming capabilities in language models. Mathematical pre-training proves advantageous for enhancing performance in language understanding and reasoning tasks.

\section{Math Pre-Training}

\subsection{Data Collection and Decontamination} This section details the methodology employed to build the DeepSeekMath Corpus utilizing data sourced from Common Crawl. Figure \ref{fig:data_collect} illustrates our stepwise pipeline, which enables the structured accumulation of a substantial mathematical corpus, beginning with a high-quality seed dataset focused on mathematics. Importantly, this iterative strategy is adaptable and may be extended to domains beyond mathematics, for instance, coding datasets.

To begin, we utilize the OpenWebMath dataset \citep{openwebmath}, a curated assemblage of high-quality mathematical web resources, as our foundational seed corpus. Leveraging this dataset, we train a fastText model \citep{joulin2016fasttext} aimed at retrieving additional web documents similar in nature to those found in OpenWebMath. For the construction of our training set, we randomly draw 500,000 samples from the seed corpus to serve as positive instances, while another 500,000 samples from Common Crawl act as negative examples. Model training is performed using an open-source implementation\footnote{\url{https://fasttext.cc}}, with the following hyperparameters: vector dimensions are set to 256, the learning rate to 0.1, the maximum word n-gram size to 3, a word occurrence threshold of 3, and a total of 3 epochs. To manage the large scale of Common Crawl, we conduct both URL-based deduplication and near-duplicate detection, which trims the dataset to 40B HTML web pages. Subsequently, the trained fastText model is used to identify mathematical web content from the deduplicated Common Crawl set. In order to eliminate mathematical pages of subpar quality, the pages are ranked according to the scores assigned by the fastText model, and only those with the highest scores are retained. We empirically determine the data volume to preserve through pre-training on the leading 40B, 80B, 120B, and 160B tokens; for the initial round, we opt to retain the top 40B tokens.

Following the initial round of data acquisition, a significant number of mathematical web pages are left out, predominantly because the fastText model had been trained on a relatively homogeneous set of positive instances. To address this limitation, we expand the seed corpus by incorporating new mathematical web sources, thereby refining the fastText model. To achieve this, we first partition the entire Common Crawl dataset into separate domains, where a domain consists of all web pages originating from the same base URL. For each such domain, we compute the proportion of web pages that were retrieved during the first collection cycle. Domains in which more than 10\% of web pages were gathered are labeled as mathematics-related (for instance, \url{mathoverflow.net}).

We then proceed to manually curate and annotate URLs that pertain to mathematical content within these math-related domains (such as \url{mathoverflow.net/questions}). Any remaining web pages associated with these annotated URLs, which were missed in prior rounds, are appended to the seed corpus. This method provides a richer set of positive examples, which, in turn, supports the training of a more effective fastText model and enhances its ability to retrieve additional mathematical content in the next cycle. Upon completing four data collection cycles, we accumulate a total of 35.5 million mathematical web pages, amounting to 120 billion tokens. By the fourth cycle, it becomes evident that nearly 98\% of the data had already been retrieved by the previous iteration, leading to the decision to terminate further data collection.

To prevent benchmark contamination, we adopt the method outlined in \cite{deepseek-coder} by excluding web pages that include either questions or answers derived from English mathematical benchmarks such as GSM8K \citep{gsm8k} and MATH \citep{MATH}, as well as from Chinese benchmarks like CMATH \citep{wei2023cmath} and AGIEval \citep{agieval}. Specifically, we remove any segment from our math training data that contains a 10-gram sequence perfectly matching a substring within the evaluation benchmarks. For benchmark materials shorter than 10 grams but at least 3 grams in length, we apply exact string matching to identify and exclude contaminated content.

\subsection{Validating the Quality of the DeepSeekMath~Corpus}

We run pre-training experiments to investigate how the DeepSeekMath~Corpus is compared with the recently released math-training corpora:

\begin{itemize}[topsep=0pt]
    \item \textbf{MathPile} \citep{mathpile}: an 8.9 billion-token, multi-source collection combining material from textbooks, Wikipedia, ProofWiki, CommonCrawl, StackExchange, and arXiv, with arXiv accounting for more than 85\% of the content;
    \item \textbf{OpenWebMath} \citep{openwebmath}: a dataset built by extracting mathematically-relevant data from CommonCrawl, resulting in a total of 13.6 billion tokens;
    \item \textbf{Proof-Pile-2} \citep{llemma}: this dataset brings together OpenWebMath, AlgebraicStack (comprised of 10.3 billion tokens of math-centric code), and 28.0 billion tokens from arXiv publications. For experiments using Proof-Pile-2, we adhere to the 2:4:1 ratio among arXiv, Web, and Code sources, following \cite{llemma}.
\end{itemize}

\subsubsection{Training Setting}
\label{sec:quality-policy}
We conduct math-centric training on a general-purpose language model that contains 1.3B parameters, constructed with the same architecture as the DeepSeek LLMs \citep{deepseek-llm}, referred to here as DeepSeek-LLM 1.3B. For each mathematical dataset, the model is independently trained over 150B tokens. All training procedures utilize the HAI-LLM \citep{haillm} framework, known for its efficiency and lightweight nature. Adhering to DeepSeek LLM training standards, we make use of the AdamW optimizer \citep{adamW} with parameters $\beta_1=0.9$, $\beta_2=0.95$, and $\mathrm{weight\_decay}=0.1$. The learning rate follows a multi-step decay policy: initially, after 2,000 warmup steps, the learning rate peaks; it is then decayed to 31.6\% of this maximum value after 80\% of training, and further dropped to 10.0\% upon completion of 90\% of training. The learning rate is capped at a maximum of 5.3e-4. For training, we use a batch size comprising 4M tokens and employ a sequence length of 4K.

\subsubsection{Evaluation Results}

\textbf{The DeepSeekMath~Corpus is of high quality, covers multilingual mathematical content, and is the largest in size.}

\begin{itemize}[topsep=0pt]
    \item \textbf{High-quality}:
    Downstream effectiveness is assessed across 8 mathematics-focused benchmarks using few-shot chain-of-thought prompting \cite{cot}. Table \ref{tab:corpora_comparison} highlights the superior results achieved by the model trained with the DeepSeekMath~Corpus. Moreover, as illustrated in Figure \ref{fig:corpora_comparisons}, this model consistently surpasses one trained on Proof-Pile-2 at the 50B token mark (equivalent to one epoch of Proof-Pile-2), which serves as evidence that the DeepSeekMath~Corpus maintains a higher average standard of quality.
    \item \textbf{Multilingual}:
    The DeepSeekMath~Corpus incorporates multilingual content, with English and Chinese being the most prevalent. Table \ref{tab:corpora_comparison} reveals that utilizing DeepSeekMath~Corpus for training leads to improved mathematical reasoning abilities in both English and Chinese. Conversely, most existing mathematical datasets, which are mainly oriented toward English, yield marginal benefits or can even impair Chinese language mathematical reasoning outcomes.
    \item \textbf{Large-scale}:
    In terms of volume, the DeepSeekMath~Corpus far exceeds the size of prior mathematical datasets. Figure \ref{fig:corpora_comparisons} demonstrates that DeepSeek-LLM 1.3B exhibits a notably steeper and more sustainable improvement trajectory when trained with the DeepSeekMath~Corpus. On the other hand, baseline datasets, which are substantially smaller and often subjected to numerous training cycles, lead to model performance that quickly stabilizes without further significant gains.
\end{itemize}

\subsection{Training and Evaluating DeepSeekMath-Base 7B}

This section presents DeepSeekMath-Base 7B, a foundational model distinguished by its advanced mathematical reasoning capabilities. Built upon DeepSeek-Coder-Base-v1.5 7B \citep{deepseek-coder}, the model underwent training on a dataset totaling 500B tokens. The dataset composition is as follows: DeepSeekMath Corpus comprises 56\%, AlgebraicStack provides 4\%, arXiv contributes 10\%, 20\% originates from Github code repositories, and the final 10\% consists of natural language text sourced from Common Crawl in both English and Chinese. The majority of training configurations align with those detailed in Section \ref{sec:quality-policy}, with notable adjustments including a peak learning rate of 4.2e-4 and a batch size set at 10M tokens.

In this work, we thoroughly evaluate the mathematical proficiency of DeepSeekMath-Base 7B. Our analysis addresses the model's capacity to generate independent, complete mathematical solutions without aid, its problem-solving effectiveness when using tools, and its competence in formal theorem proving. Additionally, we examine the broader abilities of the base model, reporting on its natural language comprehension, reasoning aptitude, and programming capabilities.

\paragraph{Mathematical Problem Solving with Step-by-Step Reasoning}

We assess the capabilities of DeepSeekMath-Base in tackling mathematical tasks through few-shot chain-of-thought prompting \citep{cot}, utilizing eight benchmarks in both English and Chinese. These benchmarks span various quantitative reasoning tasks—including GSM8K \citep{gsm8k}, MATH \citep{MATH}, and CMATH \citep{wei2023cmath}—as well as multiple-choice assessments such as MMLU-STEM \citep{mmlu} and Gaokao-MathQA \citep{agieval}. The problems addressed cover a wide range of mathematical domains and span difficulty levels from elementary school to college.

As presented in Table \ref{tab:base_math_cot}, DeepSeekMath-Base 7B demonstrates superior performance compared to other open-source base models across all eight benchmarks, including prominent models such as the popular Mistral 7B \citep{mistral} and the newly introduced Llemma 34B \citep{llemma}, which incorporates mathematical training on Proof-Pile-2 \citep{llemma}. Remarkably, DeepSeekMath-Base achieves more than a 10\% absolute improvement over other open-source base models on the challenging, competition-oriented MATH dataset, and even exceeds the performance of Minerva 540B \citep{minerva}—a significantly larger, closed-source base model (77 times the size), constructed on PaLM \citep{palm} and further specialized with mathematical content.

\paragraph{Mathematical Problem Solving with Tool Use} We assess the capability of models in mathematical problem solving on the GSM8K and MATH benchmarks by employing few-shot program-of-thought prompting as described in \citep{pot,pal}. In this approach, models are instructed to address each problem by generating Python code, with permission to make use of libraries like \textit{math} and \textit{sympy} for handling complex calculations. The resulting output from executing these programs is considered the model's answer. According to the results in Table \ref{tab:base_math_pot_proof}, DeepSeekMath-Base 7B surpasses the previous best model, Llemma 34B.

\paragraph{Formal Mathematics}

Automating formal proofs offers substantial advantages by safeguarding the correctness and dependability of mathematical proofs while also increasing productivity; as a result, this area has gained considerable momentum in recent years. In this work, we assess DeepSeekMath-Base 7B on the informal-to-formal proof translation task outlined in \citep{dsp_proof}. This task entails producing a formal proof from an informal problem statement, its corresponding formal version, and an informal solution. For our experiments, we utilize the miniF2F benchmark \citep{minif2f}, which targets the formalization of problems at Olympiad-level mathematics. Leveraging few-shot prompting, we guide the model to output formal proofs in Isabelle for each problem. In line with the methodology in \cite{dsp_proof}, we utilize the model to produce proof sketches and employ the Sledgehammer automated theorem prover \citep{sledgehammer} to fill out the details omitted by the model. Results summarized in Table \ref{tab:base_math_pot_proof} indicate that DeepSeekMath-Base 7B achieves robust results in automatic formal proof synthesis.

\paragraph{Natural Language Understanding, Reasoning, and Code}

We assess the abilities of models in natural language understanding using MMLU \citep{mmlu}, in reasoning with BBH \citep{bbh}, and in coding through the HumanEval \citep{codex} and MBPP \citep{mbpp} benchmarks. According to Table \ref{tab:understanding_reasoning_code}, DeepSeekMath-Base 7B achieves substantial gains on both MMLU and BBH when compared to the earlier DeepSeek-Coder-Base-v1.5 model \citep{deepseek-coder}, thereby underscoring the beneficial effects that mathematics-centric pretraining has on understanding and reasoning tasks. Furthermore, the integration of code tokens in continued training allows DeepSeekMath-Base 7B to retain the coding performance levels demonstrated by DeepSeek-Coder-Base-v1.5 across both HumanEval and MBPP. Taken together, DeepSeekMath-Base 7B considerably surpasses the general-purpose Mistral 7B \citep{mistral} in all three benchmarks measuring reasoning and programming competence.

\section{Supervised Fine-Tuning}

\subsection{SFT Data Curation}

Our work involves developing a comprehensive mathematical instruction-tuning dataset that spans both English and Chinese, encompassing a range of mathematical disciplines and degrees of difficulty. Each problem is matched with a corresponding solution presented in one of three formats: chain-of-thought (CoT) \citep{cot}, program-of-thought (PoT) \citep{pot,pal}, or tool-integrated reasoning \citep{tora}. The compiled dataset consists of 776K training instances.

\begin{itemize}[topsep=0pt]
    \item \textbf{English mathematical datasets}: We provide GSM8K and MATH problem sets annotated with solutions utilizing tool integration, and additionally include selected items from MathInstruct \citep{MathInstruct} and the Lila-OOD training dataset \citep{lila}, both featuring step-wise reasoning via CoT or PoT. The English datasets comprise questions from a wide spectrum of mathematical domains, including algebra, probability, number theory, calculus, and geometry.
    \item \textbf{Chinese mathematical datasets}: Our collection consists of K-12 mathematical exercises in Chinese, distributed across 76 sub-categories such as linear equations, with answers offered in both chain-of-thought and tool-augmented reasoning forms.
\end{itemize}

\subsection{Training and Evaluating DeepSeekMath-Instruct 7B}

This section presents DeepSeekMath-Instruct 7B, a model refined through mathematical instruction tuning built upon DeepSeekMath-Base. During training, examples are combined in random order until a total of 4K tokens per context is achieved. The model is optimized over 500 steps using a batch size of 256 and a fixed learning rate of 5e-5.

We evaluate models' mathematical performance both without and with tool use, on 4 quantitative reasoning benchmarks in English and Chinese. We benchmark our model against the leading models of the time:

\begin{itemize}[topsep=0pt]
    \item \textbf{Closed-source models} comprise: (1) the GPT series, with GPT-4 \citep{gpt4} and GPT-4 Code Interpreter~\footnote{\url{https://openai.com/blog/chatgpt-plugins\#\#code-interpreter}} being the most advanced, (2) Gemini Ultra and Pro \citep{gemini}, (3) Inflection-2 \citep{inflection-2}, (4) Grok-1~\footnote{\url{https://x.ai/model-card}},
    along with recently unveiled models from Chinese enterprises such as (5) Baichuan-3~\footnote{\url{https://www.baichuan-ai.com}},
    and (6) the most recent GLM-4~\footnote{\url{https://open.bigmodel.cn/dev/api\#glm-4}}

    from the GLM family \citep{glm}.

These models are designed for broad applications, with the majority having been subjected to various alignment techniques.

\item \textbf{Open-source models} encompass: general-purpose systems such as (1) DeepSeek-LLM-Chat 67B \citep{deepseek-llm}, (2) Qwen 72B \citep{qwen}, (3) SeaLLM-v2 7B \citep{seallm}, and (4) ChatGLM3 6B \citep{chatglm3}. There are also models specially adapted for mathematical tasks, including (5) InternLM2-Math 20B~\footnote{\url{https://github.com/InternLM/InternLM-Math}}, which extends InternLM2 via targeted mathematical training and subsequent instruction tuning, (6) Math-Shepherd-Mistral 7B, in which PPO training \citep{schulman2017proximal} is implemented on Mistral 7B \citep{mistral} utilizing a reward model supervised by process, (7) the WizardMath series \citep{wizardmath}, enhancing math reasoning skills in Mistral 7B and Llama-2 70B \citep{llama2} through evolve-instruct—an instruction tuning approach using AI-generated instructions—and PPO optimization, with datasets mainly from GSM8K and MATH, (8) MetaMath 70B \citep{metamath}, created by fine-tuning Llama-2 70B on an expanded dataset derived from GSM8K and MATH, (9) ToRA 34B \cite{tora}, which takes CodeLlama 34B and tunes it for mathematical reasoning incorporating tool use, and (10) MAmmoTH 70B \citep{MathInstruct}, which is an instruction-tuned version of Llama-2 70B specifically trained on MathInstruct.

Referencing Table \ref{tab:sft_rl_math}, in conditions where the use of tools is not permitted, DeepSeekMath-Instruct 7B exhibits robust capability in stepwise reasoning tasks. Specifically, when assessed on the competition-level MATH benchmark, our model outperforms all available open-source systems as well as most commercial models (such as Inflection-2 and Gemini Pro) by a margin of at least 9\% in absolute terms. This holds true even when compared to considerably larger models (e.g., Qwen 72B) or those that have undergone additional reinforcement learning targeting mathematical skills (such as WizardMath-v1.1 7B). Although DeepSeekMath-Instruct achieves comparable results to the Chinese commercial models GLM-4 and Baichuan-3 on MATH, it remains behind state-of-the-art systems like GPT-4 and Gemini Ultra.

In the evaluation scenario permitting models to combine natural language reasoning with programmatic tool utilization for problem-solving tasks, DeepSeekMath-Instruct 7B attains nearly 60\% accuracy on MATH, outperforming every other open-source model to date. Across additional benchmarks, our model demonstrates competitiveness with DeepSeek-LLM-Chat 67B—the previous leading model, despite it being an order of magnitude larger.

\subsection{Group Relative Policy Optimization} Reinforcement learning (RL) has demonstrated its utility in enhancing the mathematical reasoning skills of LLMs beyond what is achieved with Supervised Fine-Tuning (SFT) \citep{wang2023math,wizardmath}. In this section, we present Group Relative Policy Optimization (GRPO), our novel RL approach designed for efficiency and efficacy.

\subsubsection{From PPO to GRPO} Proximal Policy Optimization (PPO) \citep{schulman2017proximal} is an actor-critic RL algorithm that is widely used in the RL fine-tuning stage of LLMs \citep{ouyang2022training}. In particular, it optimizes LLMs by maximizing the following surrogate objective:

\begin{equation}
\footnotesize
    \mathcal{J}_{PPO}(\theta) = \mathbb{E}{[q \sim P(Q), o \sim \pi_{\theta_{old}}(O|q)]} \frac{1}{|o|} \sum_{t=1}^{|o|} \min \left[ \frac{\pi_\theta(o_{t} | q, o_{<t})}{\pi_{\theta_{old}}(o_{t} | q, o_{<t})} A_{t}, \text{clip} \left( \frac{\pi_\theta(o_{t} | q, o_{<t})}{\pi_{\theta_{old}}(o_{t} | q, o_{<t})}, 1 - \varepsilon, 1 + \varepsilon \right)  A_{t} \right] ,
\end{equation}

The objective function for PPO, denoted $\mathcal{J}_{PPO}(\theta)$, computes the expected value over $q$ sampled from $P(Q)$ and $o$ sampled from the policy $\pi_{\theta_{old}}(O|q)$. For each output sequence, the objective averages over all time steps $t$, evaluating the minimum of two quantities: the importance-weighted advantage and a clipped variant that restricts the ratio between the updated and the previous policy probability to within $(1 - \varepsilon, 1 + \varepsilon)$. This formulation is crucial in enforcing stable policy updates while taking into account the estimated advantage $A_t$ at each timestep.

Here, $\pi_{\theta}$ and $\pi_{\theta_{old}}$ denote the policy models corresponding to the current and previous iterations, respectively. The symbols $q$ and $o$ refer to the questions drawn from the dataset and the outputs generated by the earlier policy $\pi_{\theta_{old}}$. The hyper-parameter $\varepsilon$, specific to PPO, is utilized to regulate the clipping mechanism and thereby enhance training robustness. The term $A_t$ stands for the advantage, which is calculated through Generalized Advantage Estimation (GAE) \citep{gae}, using the subsequent rewards $\{r_{\ge t}\}$ and a value function $V_{\psi}$ that is learned concurrently. Consequently, PPO requires the concurrent optimization of a value function in addition to the policy. To prevent excessive optimization of the reward model, it is standard to incorporate a KL-divergence penalty at every token by referencing a baseline model as part of the reward calculation \citep{ouyang2022training}; that is,

\begin{equation}
    r_{t} = r_\varphi(q, o_{\le t}) - \beta \log\frac{\pi_{\theta}(o_{t}|q, o_{<t})}{\pi_{ref}(o_{t}|q, o_{<t})},
\label{eq:PPO-reward}
\end{equation}

Here, $r_\varphi$ denotes the reward model, $\pi_{ref}$ represents the reference policy—commonly set as the original SFT model—and $\beta$ serves as the scaling parameter for the KL divergence penalty.

As the value function employed in PPO is typically another model of comparable size as the policy model, it brings a substantial memory and computational burden. Additionally, during RL training, the value function is treated as a baseline in the calculation of the advantage for variance reduction. While in the LLM context, usually only the last token is assigned a reward score by the reward model, which may complicate the training of a value function that is accurate at each token. To address this, as shown in Figure \ref{fig:grpo}, we propose Group Relative Policy Optimization (GRPO), which obviates the need for additional value function approximation as in PPO, and instead uses the average reward of multiple sampled outputs, produced in response to the same question, as the baseline. More specifically, for each question $q$, GRPO samples a group of outputs $\{o_1, o_2, \cdots, o_G\}$  from the old policy  $\pi_{\theta_{old}}$  and then optimizes the policy model by maximizing the following objective:

\begin{equation}
\footnotesize
\begin{split}
    \mathcal{J}_{GRPO}(\theta) &= \mathbb{E}{[q \sim P(Q), \{o_i\}_{i=1}^G \sim \pi_{\theta_{old}}(O|q)]}  \\
    & \frac{1}{G}\sum_{i=1}^G\frac{1}{|o_i|} \sum_{t=1}^{|o_i|} \left\{ \min \left[ \frac{\pi_\theta(o_{i,t} | q, o_{i,<t})}{\pi_{\theta_{old}}(o_{i,t} | q, o_{i,<t})} \hat{A}_{i,t}, \text{clip} \left( \frac{\pi_\theta(o_{i,t} | q, o_{i,<t})}{\pi_{\theta_{old}}(o_{i,t} | q, o_{i,<t})}, 1 - \varepsilon, 1 + \varepsilon \right)  \hat{A}_{i,t} \right] - \beta \mathbb{D}_{KL}\left[\pi_{\theta} || \pi_{ref}\right]\right\} ,
\end{split}
\label{eq:GRPO-obj}
\end{equation}

where $\varepsilon$ and $\beta$ are hyper-parameters, and $\hat{A}_{i,t}$  is the advantage calculated based on relative rewards of the outputs inside each group only, which will be detailed in the following subsections. The group relative way that GRPO leverages to calculate the advantages, aligns well with the comparative nature of rewards models,  as reward models are typically trained on datasets of comparisons between outputs on the same question. Also note that, instead of adding KL penalty in the reward, GRPO regularizes by directly adding the KL divergence between the trained policy and the reference policy to the loss, avoiding complicating the calculation of  $\hat{A}_{i,t}$. And different from the KL penalty term used in (\ref{eq:PPO-reward}), we estimate the KL divergence with the following unbiased estimator \citep{kl_approx}:

\begin{equation}
\small
    \mathbb{D}_{KL}\left[\pi_{\theta} || \pi_{ref}\right] = \frac{\pi_{ref}(o_{i,t}|q,o_{i,<t})}{\pi_{\theta}(o_{i,t}|q,o_{i,<t})} - \log\frac{\pi_{ref}(o_{i,t}|q,o_{i,<t})}{\pi_{\theta}(o_{i,t}|q,o_{i,<t})} - 1,
\end{equation}

which is guaranteed to be positive.

\begin{algorithm}[t]
  \small
  \caption{Iterative Group Relative Policy Optimization}
  \textbf{Input} initial policy model $\pi_{\theta_{\text{init}}}$; reward models $r_\varphi$; task prompts $\mathcal{D}$;
  hyperparameters $\varepsilon$, $\beta$, $\mu$
  \begin{algorithmic}[1]
    \State Set current policy $\pi_\theta \leftarrow \pi_{\theta_{\text{init}}}$
    \For{iteration = 1, \dots, I}
       \State Set the reference policy $\pi_{ref} \leftarrow \pi_{\theta}$
      \For{step = 1, \dots, M}
      \State Draw a batch $\mathcal{D}_b$ from $\mathcal{D}$
      \State Store the previous policy $\pi_{\theta_{old}} \leftarrow \pi_{\theta}$
      \State For each prompt $q \in \mathcal{D}_b$, generate $G$ responses $\{o_i\}_{i=1}^G \sim \pi_{\theta_{old}} (\cdot \mid q) $
      \State Evaluate rewards $\{r_i\}_{i=1}^{G}$ for all generated outputs $o_i$ using $r_{\varphi}$
      \State Estimate group-relative advantages $\hat{A}_{i,t}$ for the $t$-th token within each $o_i$
      \For{GRPO iteration = 1, \dots, $\mu$}
        \State Optimize $\pi_{\theta}$ with respect to the GRPO criterion (Equation \ref{eq:GC-GRPO})
      \EndFor
    \EndFor \State Perform ongoing training of $r_\varphi$ utilizing a replay buffer.
    \EndFor
  \end{algorithmic}
  \textbf{Output} $\pi_\theta$
  \label{alg:iter-grpo}
\end{algorithm}

\subsubsection{Outcome Supervision RL with GRPO} Precisely, for each given question $q$, the previous policy model $\pi_{\theta_{old}}$ is utilized to generate a collection of outputs $\{o_1, o_2, \cdots, o_G\}$. Each of these outputs is evaluated by a reward model, which assigns corresponding rewards $\mathbf{r}=\{r_1, r_2, \cdots, r_G\}$ for the $G$ outputs. The assigned rewards undergo normalization: the mean reward of the group is subtracted from each reward and then the result is divided by the standard deviation of the rewards within the group. Under outcome supervision, each generated output $o_i$ receives a normalized reward at its termination, and this normalized reward is used as the advantage $\hat{A}_{i, t}$ for every token within the output; that is, $\hat{A}_{i, t} = \widetilde{r}_i = \frac{r_i- {\rm mean}(\mathbf{r})}{{\rm std}(\mathbf{r})}$. The policy is then updated to maximize the objective outlined in equation (\ref{eq:GRPO-obj}).

\subsubsection{Process Supervision RL with GRPO}  
Relying solely on outcome supervision restricts feedback to a single reward after the entire output, which can be inadequate and inefficient for guiding the policy through intricate mathematical problems. Building upon the approach in \cite{wang2023math}, we investigate process supervision, wherein feedback is provided at the conclusion of each individual reasoning step. Specifically, for a given question $q$ and $G$ sampled outputs $\{o_1, o_2, \cdots, o_G\}$, we utilize a process reward model that evaluates each step, resulting in sets of rewards: $\mathbf{R} = \{ \{r_1^{index(1)},\cdots,r_1^{index(K_1)}\}, \cdots,  \{r_G^{index(1)},\cdots,r_G^{index(K_G)}\} \}$. Here, $index(j)$ identifies the terminal token of the $j$-th reasoning step, and $K_i$ denotes the number of steps in output $i$. To ensure consistency across different outputs, these rewards are standardized using their mean and standard deviation, so that $\widetilde{r}_i^{index(j)} = \frac{r_i^{index(j)} - {\rm mean(\mathbf{R})}}{{\rm std(\mathbf{R})}}$.
Process supervision then determines the advantage for each token as the cumulative normalized reward over all subsequent steps, formally: $\hat{A}_{i, t} = \sum_{index(j) \ge t} \widetilde{r}_i^{index(j)}$.
The policy is ultimately trained by maximizing the objective described in equation (\ref{eq:GRPO-obj}).

\subsubsection{Iterative RL with GRPO} Throughout reinforcement learning training, the previously established reward model may become inadequate for effectively guiding the evolving policy model. To address this, we investigate iterative RL with GRPO. As illustrated in Algorithm \ref{alg:iter-grpo}, the iterative GRPO framework involves creating updated datasets for the reward model by sampling from the current policy model. The reward model is further refined via continual training, utilizing a replay buffer that retains 10\% of prior data. Subsequently, the current policy model is designated as the reference model, after which policy optimization continues, leveraging the updated reward model.

\subsection{Training and Evaluating DeepSeekMath-RL}

Our reinforcement learning experiments utilize the DeepSeekMath-Instruct 7B model. The RL phase draws its training data from chain-of-thought-formatted questions within the GSM8K and MATH subsets of the SFT corpus, totaling approximately 144K questions. Questions from other SFT domains are omitted to allow us to assess the effect of RL in situations where the benchmark lacks supporting data during the RL stage. Following the procedure outlined by \citep{wang2023math}, we assemble the dataset for the reward models. The initial reward model is built using DeepSeekMath-Base 7B, trained with a learning rate of $2 \times 10^{-5}$. When applying GRPO, the learning rate for the policy model is set to $1 \times 10^{-6}$, and the KL penalty coefficient is configured to $0.04$. For each question, $64$ completions are sampled. Both the maximum sequence length and the batch size are fixed at $1024$. The policy model undergoes only one update after every exploration stage. Evaluation of DeepSeekMath-RL 7B on benchmark datasets mirrors the approach used for DeepSeekMath-Instruct 7B. For DeepSeekMath-RL 7B, the GSM8K and MATH tasks with chain-of-thought reasoning serve as in-domain benchmarks, while all remaining datasets are categorized as out-of-domain.

Table \ref{tab:sft_rl_math} presents the results for both open- and closed-source models equipped with chain-of-thought and tool-integrated reasoning across English and Chinese evaluation tasks. Our observations are as follows: 1) When employing chain-of-thought reasoning, DeepSeekMath-RL 7B achieves accuracy scores of 88.2\% on GSM8K and 51.7\% on MATH, outperforming every open-source model within the 7B–70B scale, and even surpassing most closed-source alternatives. 2) Notably, DeepSeekMath-RL 7B’s training exclusively involved instruction tuning on chain-of-thought data from GSM8K and MATH, building upon DeepSeekMath-Instruct 7B. Although its training dataset is limited, DeepSeekMath-RL 7B consistently exceeds the performance of DeepSeekMath-Instruct 7B on all assessment criteria, illustrating the substantial impact of reinforcement learning.

\section{Discussion}
Here, we present the outcomes observed in both our pre-training and reinforcement learning (RL) experimental setups.

\subsection{Lessons Learnt in Pre-Training}

We begin by detailing our observations from the pre-training phase. Except where noted, all experiments follow the training configurations described in Section \ref{sec:quality-policy}. Additionally, throughout this section, references to the DeepSeekMath Corpus pertain to an 89B-token dataset derived from the second round of data collection.

\subsubsection{Code Training Benefits Mathematical Reasoning}

A widely discussed but unproven theory posits that exposure to code enhances reasoning capabilities. In this work, we seek to partially address this claim in the context of mathematics: code training strengthens models' mathematical reasoning skills regardless of whether they utilize tools.

To study how code training affects mathematical reasoning, we experimented with the following two-stage training and one-stage training settings:

\textbf{Two-Stage Training}

\begin{itemize}[topsep=0pt]
    \item \textbf{Code Training for 400B Tokens $\rightarrow$ Math Training for 150B Tokens}: In this configuration, DeepSeek-LLM 1.3B is initially pretrained on a dataset of 400B code tokens, which is then succeeded by additional training on 150B math-specific tokens;
    \item \textbf{General Training for 400B Tokens $\rightarrow$ Math Training for 150B Tokens}: To provide a baseline, we conduct a parallel experiment where, in the initial training phase, general-purpose tokens (drawn from an extensive general text corpus compiled by DeepSeek-AI) are used in place of code tokens. This aims to assess whether utilizing code tokens confers any benefit for enhancing the model's mathematical reasoning, relative to general tokens.
\end{itemize}

\textbf{One-Stage Training}

\begin{itemize}[topsep=0pt]
    \item \textbf{150B Math Token Training}: DeepSeek-LLM 1.3B is trained using 150B tokens specifically from mathematical datasets;
    \item \textbf{Combined Training with 400B Code Tokens and 150B Math Tokens}: Sequential math training after initial code training has been observed to reduce proficiency on coding tasks. To address this, we explore simultaneous training with a blend of code and math tokens, assessing if such integration can enhance mathematical reasoning capabilities and mitigate the effects of catastrophic forgetting.
\end{itemize}

\paragraph{Results}
The downstream performance across various training configurations is presented in Table \ref{tab:code-math} and Table \ref{tab:code-math-general-eval}.

Incorporating code training enhances program-aided mathematical reasoning across both two-stage and one-stage training paradigms. Table \ref{tab:code-math} illustrates that, in the context of two-stage training, applying code training exclusively already results in notable gains on GSM8K and MATH problem-solving tasks that require Python solutions. Introducing math training in the subsequent phase produces additional improvements. Notably, for one-stage training, jointly including code tokens with math tokens not only addresses the catastrophic forgetting issue encountered in two-stage training but also promotes synergy between coding capabilities (Table \ref{tab:code-math-general-eval}) and program-aided mathematical reasoning (Table \ref{tab:code-math}).

Code training contributes to enhancements in mathematical reasoning even when external tools are not utilized. When adopting a two-stage training framework, introducing code training in the first phase leads to measurable gains early on. This approach further increases the effectiveness of subsequent math-focused training, culminating in optimal performance. In contrast, a single-stage training process that blends code tokens and math tokens appears to diminish mathematical reasoning in the absence of tool support. A possible explanation is that DeepSeek-LLM 1.3B, given its relatively small model size, may not possess sufficient capacity to concurrently integrate both code and math information effectively.

\subsubsection{ArXiv Papers Seem Ineffective in Improving Mathematical Reasoning}

ArXiv papers are commonly included as a component of math pre-training data \citep{minerva,gpt-f,llemma,mathpile}. However, detailed analysis regarding their impact on mathematical reasoning has not been extensively conducted. Perhaps counter-intuitively, according to our experiments, arXiv papers seem ineffective in improving mathematical reasoning. We experiment with models of different sizes, including DeepSeek-LLM 1.3B and DeepSeek-Coder-Base-v1.5 7B \citep{deepseek-coder}, using arXiv corpora that underwent varied processing pipelines:

\begin{itemize}[topsep=0pt]
    \item \textbf{MathPile} \citep{mathpile}: comprises 8.9 billion tokens, curated using specialized cleaning and filtering approaches, with more than 85\% sourced from scientific papers on arXiv;
    \item \textbf{ArXiv-RedPajama} \citep{redpajama}: consists of all arXiv LaTeX documents, systematically stripped of preambles, comments, macros, and bibliographies, resulting in a dataset of 28.0 billion tokens.
\end{itemize}

During our study, we independently trained DeepSeek-LLM 1.3B on 150B tokens and DeepSeek-Coder-Base-v1.5 7B on 40B tokens, utilizing distinct arXiv corpora. Our findings indicate that incorporating arXiv papers does not substantially enhance mathematical reasoning capabilities. When exposed solely to arXiv content, neither model demonstrates marked progress; in fact, performance stagnates or declines across a range of mathematical evaluation tasks with varying difficulty, as used in our analysis. These assessments comprise quantitative reasoning datasets such as GSM8K and MATH (Table \ref{tab:arxiv-cot}), multiple-choice exams including MMLU-STEM (Table \ref{tab:arxiv-cot}), as well as formal mathematics tasks like miniF2F (Table \ref{tab:arxiv_atp}).

However, this conclusion has its limitations and should be taken with a grain of salt. We have not yet studied:

\begin{itemize}[topsep=0pt]
    \item How arXiv tokens influence additional mathematical tasks not covered by the current study, for example, transforming formal theorems or proofs into their informal representations;
    \item The outcomes arising from integrating arXiv tokens with alternative datasets;
    \item If advantages derived from using arXiv papers persist or increase as model size grows.
\end{itemize}

Thus, further exploration is required, which we leave for future studies.

\subsection{Insights of Reinforcement Learning}
\subsubsection{Toward a Unified Paradigm} This section introduces a comprehensive framework to systematically examine various training approaches, including SFT, RFT, DPO, PPO, GRPO. Additionally, we perform experimental investigations to identify and analyze the key components of this unified paradigm. Typically, for any given training algorithm, the gradient with respect to the parameters $\theta$ is expressed as follows:

\begin{equation}
    \nabla_{\theta}\mathcal{J}_{\textcolor{red}{\mathcal{A}}}(\theta) = \mathbb{E}[\underbrace{(q,o) \sim \textcolor{red}{\mathcal{D}}}_{Data \ Source}]\left( \frac{1}{|o|} \sum_{t=1}^{|o|}  \underbrace{GC_{{\mathcal{A}}}(q, o, t, \textcolor{red}{\pi_{{rf}}})}_{Gradient \ Coefficient}  \nabla_{\theta}\log \pi_{\theta}(o_t | q, o_{<t})\right).
\label{eq:objective}
\end{equation}

There exist three key components: 1) \textit{Data Source $\mathcal{D}$}, which determines the training data; 2) \textit{Reward Function $\pi_{{rf}}$}, which is the source of the training reward signal; 3) \textit{Algorithm $\mathcal{A}$}: which processes the training data and the reward signal to the gradient coefficient $GC$ that determines the magnitude of the penalty or reinforcement for the data. We analyze several representative methods based on such a unified paradigm:

\begin{itemize}
    \item \textbf{Supervised Fine-tuning (SFT)}: In SFT, the pretrained model undergoes further training using a dataset curated and labeled by humans.
    \item \textbf{Rejection Sampling Fine-tuning (RFT)}: RFT extends the SFT procedure by additionally fine-tuning on those outputs, generated in response to SFT data, that are selectively retained according to the accuracy of their answers.
    \item \textbf{Direct Preference Optimization (DPO)}: DPO builds on the SFT foundation by applying fine-tuning to the model with enhanced outputs created from the SFT model, leveraging a pairwise DPO loss to optimize preferences.
    \item \textbf{Online Rejection Sampling Fine-tuning (Online RFT)}: Unlike traditional RFT, Online RFT starts with the SFT-initialized policy model and continuously updates it through fine-tuning on freshly generated outputs sampled directly from the current policy during training.
    \item \textbf{PPO/GRPO}: With PPO/GRPO, the approach is to initialize the policy with the SFT model and subsequently enhance it using online outputs sampled during the reinforcement learning process with the live policy model.
\end{itemize}

Table \ref{tab:data_policy} provides an overview of the elements comprising these methods. For a comprehensive explanation of the derivation, see Appendix \ref{app:analysis-rl}.

\paragraph{Observation about Data Source} We classify data sources into two main types: online sampling and offline sampling. In online sampling, the training data originates from the outputs produced by the policy model as it explores in real time during training. Conversely, offline sampling utilizes data generated by the initial SFT model as the training samples. Methods such as RFT and DPO employ the offline approach, whereas Online RFT and GRPO make use of the online approach.

As illustrated in Figure \ref{fig:rl-analysis}, our results indicate that Online RFT achieves notably better performance than RFT across both benchmarks. At the start of training, Online RFT and RFT yield similar results; however, as training progresses, Online RFT attains a marked performance lead, highlighting the benefits of online learning. This pattern aligns with expectations: early on, the actor remains similar to the SFT model, resulting in sampled data with only slight variations. As training advances, the divergence in data generated by the actor increases, and the advantages of on-the-fly data sampling become more apparent.

\paragraph{Observation about Gradient Coefficient} The method translates the reward signal into a gradient coefficient, which is then used to update the model parameters. In our experimental setup, the reward function is categorized into `Rule' and `Model'. `Rule' entails evaluating the appropriateness of a response purely on answer correctness, whereas `Model' indicates that a separate reward model is trained to assign scores to responses. The reward model relies on training data derived from rule-based assessments. As shown in Equations \ref{eq:GC-RFT} and \ref{eq:GC-GRPO}, there exists an important distinction between GRPO and Online RFT: GRPO dynamically modifies the gradient coefficient as a function of the reward model's output. This enables GRPO to administer varying levels of encouragement or discouragement to responses, depending on the exact reward magnitude. In comparison, Online RFT does not possess this flexibility; it neither penalizes erroneous responses nor differentiates among correct ones, consistently providing reinforcement of equal strength to every correct response.

As illustrated in Figure \ref{fig:rl-analysis}, GRPO outperforms online RFT, underscoring the effectiveness of adjusting the positive and negative gradient coefficients. Moreover, GRPO+PS achieves better results than GRPO+OS, suggesting that incorporating step-aware, fine-grained gradient coefficients offers additional advantages. We also investigate iterative RL by performing two successive iterations in our experiments. According to Figure \ref{fig:iter-rl}, iterative RL leads to marked enhancements in performance, with the most pronounced gains occurring during the initial iteration.

\subsubsection{Why RL Works?} In our study, reinforcement learning is applied using a selective set from the instruction tuning dataset, leading to notable gains over models solely relying on instruction tuning. To clarify the underlying mechanism for RL's effectiveness, we assess both Pass@K and Maj@K metrics for Instruct and RL-based models across two evaluation benchmarks. As presented in Figure \ref{fig:maj-pass}, reinforcement learning specifically increases Maj@K scores, while Pass@K remains unchanged. This suggests that RL’s contribution lies in strengthening the reliability of the model's output distribution, indicating that \textbf{the observed advancements mainly stem from elevating the prominence of correct answers among the TopK choices, rather than from a fundamental improvement in the core modeling capability.} Likewise, \citep{wang2023making} report a \textbf{misalignment problem} in reasoning scenarios when using SFT models, and demonstrate that integrating various preference alignment techniques can boost reasoning performance in these models \citep{yuan2023rrhf,song2023preference,wang2023making}.

\subsubsection{How to Achieve More Effective RL?}
\label{sec:effec_rl}
We show that RL performs effectively on mathematical reasoning challenges. Furthermore, we present a general framework that connects various prominent training strategies, viewing all of them as forms of either direct or approximated RL approaches. This framework, as outlined in Equation \ref{eq:objective}, centers on three essential elements: Data Source, Algorithm, and Reward Function. We discuss several prospective avenues for further exploration regarding these three factors.

\paragraph{Data Source} The data source forms the foundational input for all training approaches. Within reinforcement learning (RL), this refers specifically to unlabeled questions paired with the responses generated by the policy model. Throughout this study, we restrict ourselves to using only the questions provided during the instruction tuning phase, generating responses via a straightforward nucleus sampling approach. We hypothesize that this limitation might account for the observed RL pipeline improvements being confined primarily to Maj@K metrics. In subsequent work, we aim to investigate how our RL framework performs with question prompts that are out-of-distribution, together with \textbf{advanced sampling (decoding) strategies}, such as those utilizing tree-search methods \citep{yao2023tree}. Additionally, \textbf{efficient inference techniques} \citep{xia-etal-2023-speculative,leviathan2023fast,kwon2023efficient,xia2024unlocking}, which are crucial for enhancing the exploration capabilities of policy models, also represent a key aspect influencing performance.

\paragraph{Algorithms} Algorithms utilize both the input data and the corresponding reward signal in calculating the gradient coefficient, which in turn updates the parameters of the model. As suggested in Equation \ref{eq:objective}, prevailing approaches effectively place full \textbf{TRUST} in the feedback provided by the reward function to adjust—either upward or downward—the conditional likelihood of specific tokens. Nevertheless, relying exclusively on the reward signal assumes its reliability, which cannot always be guaranteed, particularly when addressing highly complex tasks. Illustratively, even in the case of PRM800K datasets \citep{lightman2023let}, for which annotation was performed by skilled experts, the rate of incorrect labeling remains around 20\%\footnote{\url{https://github.com/openai/prm800k/issues/12\#issuecomment-1728491852}}. Therefore, our work investigates reinforcement learning methods that demonstrate robustness in the presence of noisy reward signals. We posit that alignment techniques inspired by \textbf{WEAK-TO-STRONG} \citep{burns2023weak} have the potential to fundamentally transform current learning algorithms.

\paragraph{Reward Function} The reward function serves as the foundational signal guiding training. In reinforcement learning (RL), this function is typically embodied by a neural reward model. Three pivotal research directions for reward models can be identified: 1) \textbf{Enhancing the generalization capacity of the reward model.} It is crucial that the reward model generalizes well, especially when faced with out-of-distribution queries or complex outputs from advanced decoders; failing to do so risks RL merely maintaining the status quo of large language model (LLM) distributions rather than driving substantive improvement; 2) \textbf{Capturing and expressing uncertainty in the reward model.} Properly accounting for uncertainty can provide a vital connection between less robust reward models and algorithms for weak-to-strong learning; 3) \textbf{Developing high-quality process reward models with greater efficiency,} such that they can deliver granular, process-level signals to facilitate reasoning during training \citep{lightman2023let,wang2023math}.

\section{Conclusion, Limitation, and Future Work}

In this work, we introduce DeepSeekMath, which surpasses all existing open-source models on the MATH benchmark at the competition level and achieves performance comparable to that of closed-source models. DeepSeekMath~is built upon DeepSeek-Coder-v1.5 7B, followed by continual training using 500B tokens, where 120B math-focused tokens are derived from Common Crawl data. Through thorough ablation studies, we demonstrate that web-sourced materials offer considerable value in constructing high-quality mathematical datasets, whereas data from arXiv do not yield the anticipated benefits. We propose Group Relative Policy Optimization (GRPO), a modified form of Proximal Policy Optimization (PPO), capable of substantially enhancing mathematical reasoning with reduced memory requirements. Experimental evidence indicates GRPO's efficacy remains robust even for DeepSeekMath-Instruct 7B at advanced benchmark scores. Additionally, we present a unified framework for conceptualizing various techniques and highlight several promising avenues to further improve reinforcement learning approaches.

While DeepSeekMath~demonstrates strong performance on benchmarks assessing quantitative reasoning, its proficiency in geometry and theorem proving falls short compared to proprietary models. For example, during our preliminary assessment, the model was unable to solve questions concerning triangles and ellipses, which suggests there may be a bias stemming from the choice of data during pre-training and fine-tuning. Furthermore, the limited scale of the model constrains DeepSeekMath's~few-shot learning abilities, rendering it less effective than GPT-4 in this respect; unlike GPT-4, which shows noticeable gains with few-shot input, DeepSeekMath~performs nearly identically on both zero-shot and few-shot tasks. Going forward, we plan to refine our data engineering pipeline in order to build a more robust and diverse pre-training dataset. Additionally, we aim to investigate novel approaches to reinforcement learning for LLMs, as discussed in Section \ref{sec:effec_rl}.

\end{document}